\section{Finite translation orbits on the ground space}
\label{sec:hastings_finite_translation_orbit}

The argument at\cite[lines 1644--1654]{Hastings2006Solving} concerns
translations whose momentum denominators may grow with the system size.
The finite-dimensional calculation below replaces a fixed supercell by
finitely many translates. Its constants depend on the dimension of the
ground space, without a condition on its translation eigenvalues.

\begin{definition}[Translation on the ground space]
    \label{def:hastings_ground_translation}
    \lean{TNLean.Lattice.groundTranslation,
        TNLean.Lattice.groundTranslationConjugation}
    \leanok
    Let $H$ be invariant under unitary site translations $U_g$, and let
    $K=\ker(H-E_0)$ be its ground eigenspace. Write $u_g=U_g|_K$.
    The induced translation on the ground-space operator algebra is
    $T_g(B)=u_gBu_g^{-1}$.
\end{definition}

\begin{lemma}[Compression and translation]
    \label{lem:hastings_ground_translation_compression}
    \lean{Matrix.sitePermMatrix_one, TNLean.Lattice.sitePerm_starProjection_commute,
        TNLean.Lattice.groundTranslation_apply_coe,
        TNLean.Lattice.groundTranslation_symm_apply,
        TNLean.Lattice.groundTranslationConjugation_groundSpaceCompression,
        TNLean.Lattice.groundTranslationConjugation_pow_compression_of_covariant,
        TNLean.Lattice.groundTranslationConjugation_pow_compression}
    \leanok
    Let $P$ be the orthogonal projection onto $K$, and put
    $C_K(W)=(PWP)|_K$. Then $U_gP=PU_g$, $u_g^{-1}=u_{g^{-1}}$, and
    \[
        T_g(C_K(W))=C_K(U_gWU_g^*).
    \]
    Consequently $T_g^k(C_K(W))=C_K(U_{g^k}WU_{g^k}^*)$ for every $k\ge0$.
    If $M_{g\cdot i}=U_gM_iU_g^*$, this also gives
    $T_g^k(C_K(M_i))=C_K(M_{g^k\cdot i})$.
\end{lemma}
\begin{proof}\leanok
    \uses{def:hastings_ground_translation}
    Commutation with $H$ preserves $K$. The inverse translation also
    preserves $K$, so its image is exactly $K$. An isometry transports
    orthogonal projections onto the corresponding image subspaces;
    hence $U_gP=PU_g$. Restriction gives the inverse and conjugation
    identities, and induction gives the identity for powers.
\end{proof}

\begin{lemma}[Norms after restriction to the ground space]
    \label{lem:hastings_ground_restriction_norm}
    \lean{TNLean.Lattice.norm_groundSpaceCompression_eq,
        TNLean.Lattice.norm_groundSpaceCompression_commutator_eq}
    \leanok
    For all operators $W,B$,
    \[
        \|C_K(W)\|=\|PWP\|,\qquad
        \|[C_K(W),C_K(B)]\|=\|[PWP,PBP]\|.
    \]
\end{lemma}
\begin{proof}\leanok
    Each ambient operator in these identities vanishes on $K^\perp$
    and has its range in $K$. Its norm is therefore the norm of its
    restriction to $K$; equivalently, apply it to $Px$ and use
    $\|Px\|\le\|x\|$.
\end{proof}

\begin{lemma}[A finite family of separated translates]
    \label{lem:hastings_finite_translation_spacing}
    \lean{TNLean.Lattice.coordinate_pow_smul,
        TNLean.Lattice.natAbs_valMinAbs_add_mul_ge_of_near,
        TNLean.Lattice.edist_pow_mul_ge_of_near}
    \leanok
    Let $\pi$ be a cell coordinate modulo $N$, and suppose translation
    $e$ increases $\pi$ by one. Fix integers $L>0$ and $m\ge0$ with
    $2(m+1)L\le N$. If the circular coordinate separation of $i,j$ is
    at most $\lfloor L/2\rfloor$, then the separation of
    $i,e^{kL}\cdot j$ is at least $\lfloor L/2\rfloor$ for every
    $1\le k\le m$. If each graph edge changes the coordinate by at
    most one, their graph distance satisfies the same lower bound.
\end{lemma}
\begin{proof}\leanok
    Translation by $e^r$ adds $r$ to the coordinate, by induction.
    Represent the initial difference by a signed residue of absolute
    value at most $\lfloor L/2\rfloor$. Adding $kL$ places it between
    $\lfloor L/2\rfloor$ and $N-\lfloor L/2\rfloor$, by the assumed
    bound on $mL$. This proves the circular separation, including a
    residue crossing zero. Summing the coordinate changes along a
    graph path gives the graph-distance bound.
\end{proof}

\begin{theorem}[A commutator from finitely many translation iterates]
    \label{thm:hastings_finite_translation_commutator}
    \lean{TNLean.Lattice.norm_commutator_le_of_finite_translation_orbit_bounds}
    \leanok
    Let $K$ have complex dimension $n$, and let $T$ be a complex linear
    isometry of its operator space. If $C\ge0$ and
    \[
        \|[A,T^k(B)]\|\le C\qquad(1\le k\le n^2),
    \]
    then $\|[A,B]\|\le(2^{n^2}-1)C$.
\end{theorem}
\begin{proof}\leanok
    Put $m=n^2$ and let $p(z)=\sum_{k=0}^m p_kz^k$ be the
    characteristic polynomial of $T$. A nonzero eigenvector and norm
    preservation imply that every root has modulus one. Thus
    $|p_0|=1$, and the coefficient bound for a monic polynomial with
    roots on the unit circle gives
    $\sum_{k=1}^m|p_k|\le2^m-1$.
    Cayley--Hamilton gives
    \[
        B=-\sum_{k=1}^m p_0^{-1}p_k T^k(B).
    \]
    Apply the linear map $B\mapsto[A,B]$ and the triangle inequality.
    The operator space has dimension $n^2$, as used above. When
    $n=0$, all its operators vanish and the conclusion still holds.
\end{proof}

\begin{lemma}[The integer orbit spacing]
    \label{lem:hastings_translation_floor_decay}
    \lean{TNLean.Lattice.translation_half_natDiv_lower_bound,
        TNLean.Lattice.exp_neg_half_translation_natDiv_le,
        TNLean.Lattice.translation_finiteOrbit_rate_pos}
    \leanok
    For natural numbers $N,m$, put $L=\lfloor N/(2(m+1))\rfloor$.
    Then
    \begin{align}
        \frac{N}{4(m+1)} & \le\frac L2+\frac12, \\
        e^{-\gamma L/2}  & \le e^{\gamma/2}
        e^{-\gamma N/(4(m+1))}\qquad(\gamma\ge0).
        \label{eq:hastings_translation_floor_decay}
    \end{align}
    If $\gamma>0$, the rate $\gamma/(4(m+1))$ is positive.
\end{lemma}
\begin{proof}\leanok
    Euclidean division gives $N=2(m+1)L+r$ with $0\le r<2(m+1)$.
    Divide by $4(m+1)$ and exponentiate the resulting inequality.
\end{proof}

\begin{lemma}[The support offset in the translation estimate]
    \label{lem:hastings_translation_support_offset_decay}
    \lean{TNLean.Lattice.translation_exp_floor_decay_le}
    \leanok
    \uses{lem:hastings_translation_floor_decay}
    With $L=\lfloor N/(2(m+1))\rfloor$, $\gamma\ge0$ and $\rho\in\mathbb R$,
    \begin{align}
        e^{-\gamma(L/2-2\rho-1)}
        \le e^{\gamma(2\rho+2)}e^{-\gamma N/(4(m+1))}.
        \label{eq:hastings_translation_support_offset_decay}
    \end{align}
\end{lemma}
\begin{proof}\leanok
    \uses{lem:hastings_translation_floor_decay}
    Factor the left side as $e^{\gamma(2\rho+1)}e^{-\gamma L/2}$.
    The rounding loss in
    (\ref{eq:hastings_translation_floor_decay}) gives the prefactor
    $e^{\gamma(2\rho+3/2)}$, which is at most $e^{\gamma(2\rho+2)}$.
\end{proof}

\begin{lemma}[A common estimate for small and large systems]
    \label{lem:hastings_translation_near_far_decay}
    \lean{TNLean.Lattice.one_le_two_pow_sub_one,
        TNLean.Lattice.translation_finiteOrbit_decay_of_near_far}
    \leanok
    \uses{lem:hastings_translation_floor_decay,
        lem:hastings_translation_support_offset_decay}
    Let $m\ge1$, $\gamma>0$, $C\ge0$, $\rho,z\in\mathbb R$, and set
    $L=\lfloor N/(2(m+1))\rfloor$. Suppose $z\le C$ and, whenever
    $L/2>2\rho+1$, suppose also
    $z\le(2^m-1)C e^{-\gamma(L/2-2\rho-1)}$. Then for every $N$,
    \begin{align}
        z\le(2^m-1)C e^{\gamma(2\rho+2)}e^{-\gamma N/(4(m+1))}.
        \label{eq:hastings_translation_near_far_decay}
    \end{align}
    This includes $L=0$ and all other small sizes with the same constants.
\end{lemma}
\begin{proof}\leanok
    \uses{lem:hastings_translation_floor_decay,
        lem:hastings_translation_support_offset_decay}
    The coefficient satisfies $2^m-1\ge1$. For $L/2>2\rho+1$, apply
    (\ref{eq:hastings_translation_support_offset_decay}) to the far estimate.
    Otherwise (\ref{eq:hastings_translation_floor_decay}) gives
    $N/(4(m+1))\le2\rho+2$, so
    $e^{\gamma(2\rho+2)}e^{-\gamma N/(4(m+1))}\ge1$ and the bound follows
    from $z\le C$.
\end{proof}

\begin{lemma}[Translation of one observable]
    \label{lem:hastings_finite_observable_orbit}
    \lean{TNLean.Lattice.norm_groundCompression_commutator_le_of_finite_observable_orbit}
    \leanok
    Let $H$ be translation invariant, let $m\ge\max\{1,(\dim K)^2\}$,
    and suppose $2(m+1)L\le N$. If every translate of $B$ whose centre
    is at graph distance at least $\lfloor L/2\rfloor$ from the centre
    of $W$ satisfies $\|[PWP,PU_gBU_g^*P]\|\le C$, where $C\ge0$,
    then $\|[PWP,PBP]\|\le(2^m-1)C$.
\end{lemma}
\begin{proof}\leanok
    \uses{lem:hastings_finite_translation_spacing,
        thm:hastings_finite_translation_commutator,
        lem:hastings_ground_translation_compression,
        lem:hastings_ground_restriction_norm}
    For initially separated centres use the identity translation.
    Otherwise translate by $e^{kL}$, for $1\le k\le(\dim K)^2$,
    and apply the finite-dimensional recurrence to $T_{e^L}$.
    Norms before and after restriction to $K$ agree. If $L=0$, the
    original pair already satisfies the hypothesis. Finally
    $1\le2^m-1$ absorbs the cases needing no nontrivial translate.
\end{proof}

\begin{theorem}[Exponential smallness under invariant-Hamiltonian translations]
    \label{thm:hastings_finite_translation_exp}
    \lean{TNLean.Lattice.norm_groundCompression_commutator_le_finite_translation_exp}
    \leanok
    Suppose translations preserve the interaction graph and the
    Hamiltonian $H$. Let a cell coordinate modulo $N$ change by at
    most one along every graph edge, and let $e$ advance it by one.
    Assume the local generator bounds of
    Theorem~\ref{thm:hastings05_growth_free_blecm_separation}, and let
    the original spectral gap above $K$ be $\Delta>0$.
    If $\dim K\le n_0$, put
    \[
        m=\max\{1,n_0^2\},\qquad
        v=4v_R|b_0|e^{2R}+1,\qquad
        \gamma=\min\{1/4,\Delta/(8v)\}>0.
    \]
    Let $W,B$ act within graph radius $\rho$ of $i,j$, respectively,
    and let the first support have at most $v_{\rm obs}$ sites. Then
    \begin{align}
        \|[PWP,PBP]\|
         & \le (2^m-1)(4v_{\rm obs}+4)\|W\|\|B\|
        e^{\gamma(2\rho+2)}e^{-\gamma N/[4(m+1)]}.
        \label{eq:hastings_finite_translation_exp}
    \end{align}
    There is no condition on the ground-state momenta or on covariance
    of a chosen decomposition of $H$.
\end{theorem}
\begin{proof}\leanok
    \uses{lem:hastings_finite_observable_orbit,
        thm:hastings05_growth_free_blecm_separation,
        lem:hastings_translation_near_far_decay}
    Set $L=\lfloor N/[2(m+1)]\rfloor$. A translation preserves the
    second support's radius and its operator norm. If its centre is
    at distance at least $D$ from $i$, the support separation is at
    least $D-2\rho$. The propagation and original-gap estimate bounds
    the commutator by
    $(4v_{\rm obs}+4)\|W\|\|B\|e^{-\gamma(D-2\rho)}$, with natural
    subtraction when the separation is an integer. Apply the
    preceding lemma at $D=\lfloor L/2\rfloor$. The floor and support
    losses are absorbed by the scalar estimates, yielding
    (\ref{eq:hastings_finite_translation_exp}). For small $N$, use
    $\|[PWP,PBP]\|\le2\|W\|\|B\|$ and the same prefactor.
\end{proof}

\begin{corollary}[The actual truncated Gaussian terms without momentum restrictions]
    \label{cor:hastings_finite_translation_truncated}
    \lean{TNLean.Lattice.norm_groundCompression_truncated_commutator_le_finite_translation_exp}
    \leanok
    In the preceding setting let $M_i$ be the actual Gaussian filter
    of $h_i$ under the Hamiltonian retaining a prescribed set of
    interactions near $i$. Suppose the initial term and every retained
    interaction act within graph radius $\rho$ of $i$, in a set of
    at most $v_{\rm obs}$ sites, and $\|h_i\|\le J$ for all $i$.
    Then
    \[
        \|[PM_iP,PM_jP]\|
        \le (2^m-1)(4v_{\rm obs}+4)J^2
        e^{\gamma(2\rho+2)}e^{-\gamma N/[4(m+1)]}.
    \]
    The retained sets need not be covariant. This includes the
    physical ball cutoff of the source and the case of momentum
    denominators growing with $N$.
\end{corollary}
\begin{proof}\leanok
    \uses{thm:hastings_finite_translation_exp,
        thm:hastings_truncated_support,thm:hastings_gaussian_filtered_spectral}
    The retained Hamiltonian and the initial term have the stated
    support, so their Gaussian evolution and its integral have that
    support. Gaussian filtering contracts the norm, giving
    $\|M_i\|\le J$. Apply the theorem with $W=M_i$, $B=M_j$.
\end{proof}
